% MAS1004 Assignment 1: long report.
%
% You fill this in with your agent, in your repository. AGENTS.md tells the
% agent how. Red text in the PDF is what still has to be replaced.
%
% - Pictures go in figures/. Copy them from results/, which is not in git.
% - Outputs go in outputs/. run.py saves what it prints in
%   results/<tag>_output.txt, and PLAN.md saves the other outputs in results/
%   with `| tee results/...`. Copy them here. \outputfile{outputs/check.txt}
%   prints a file exactly as it is. Every number in this report has to come
%   from such an output.
%
% To make the PDF: compress this folder into a zip, and on Overleaf choose
% New Project, Upload Project. Set Menu, Compiler to XeLaTeX, and Recompile.

\documentclass[11pt]{article}
\usepackage{mas1004}

\reporttitle{Assignment 1: Long report}
\studentname{Your name}
\studentid{Your student ID}

\begin{document}
\makeheader

% ---------------------------------------------------------------------------
\section*{Problem 1. Categories and images}

\subsection*{Why these categories}
I had just seen the new Chiikawa film, so I wanted to work with the characters
of that comic rather than with cups or flowers. I collected photographs of the
plush toys and goods of three of its characters instead of much wider
categories, because the three look alike and I wanted to see whether the model
can tell them apart at all, and what it uses when it does.

\subsection*{What I expected the model to use}
Before training I expected it to use the whole outline: the shape of the eyes
and the eyebrows, the ears, the tail, and the colour the character is known by.
I did not expect the background to matter, and I expected Hachiware and Usagi,
the two with colour on them, to be easier than Chiikawa, which is plain white.

\subsection*{Images per category}
150 images for each of the three categories, 450 in all, as downloaded; the
counts are in the first lines of \texttt{results/run\_output.txt}. After
cleaning, \texttt{data/clean} holds 104 chiikawa, 110 hachiware and 99 usagi
images (\texttt{check.py}, and the table in Problem 4).

% ---------------------------------------------------------------------------
\section*{Problem 2. Preparing an image}

\sidebyside{figures/prepare_original.jpg}{hachiware\_character/0093.jpg as
downloaded: 1080x1920 on screen, resized to 256x455 here, with a red box around
the 224x224 that \texttt{prepare\_image} keeps}
           {figures/prepare_result.png}{what \texttt{prepare\_image} made of it:
224x224, the only thing the model ever sees}

\texttt{prepare\_image} resizes the shorter side to 256 and then keeps only the
middle 224x224 square, so on a tall picture it throws the top and the bottom
away. In hachiware\_character/0093.jpg the bow and the top of the head are
inside the square and everything below them is gone: the eyes, the mouth and the
pink cheeks, which are the only things that tell Hachiware from Chiikawa, never
reach the model. Every 224x224 value is then divided by 255 and shifted by the
ImageNet mean and standard deviation, channel by channel, with the channels
first, so the model reads the picture exactly the way it read the million
ImageNet photographs. This image is still in \texttt{data/clean}: the model
trains on a bow.

% ---------------------------------------------------------------------------
\section*{Problem 3. First training runs}

\begin{center}\footnotesize
\begin{tabular}{llllrrrr}
\toprule
tag & start & trained & epochs & lr & trainable parameters & train acc. & test acc. \\
\midrule
run     & ImageNet & all layers & 10 & 0.0001 & 11,178,051 & 96.1\% & 83.3\% \\
scratch & random   & all layers & 10 & 0.0001 & 11,178,051 & 95.8\% & 66.7\% \\
frozen  & ImageNet & last layer & 10 & 0.001  & 1,539      & 85.8\% & 74.4\% \\
long    & ImageNet & all layers & 20 & 0.0001 & 11,178,051 & 95.6\% & 84.4\% \\
\bottomrule
\end{tabular}
\end{center}

\picturefile{figures/run_curves.png}

\sidebyside{figures/scratch_curves.png}{scratch: random weights, 10 epochs, the
same pictures and the same code as run}
           {figures/frozen_curves.png}{frozen: only the last layer is trained,
with the larger learning rate}

\picturefile[0.6\linewidth]{figures/long_curves.png}

With 20 epochs instead of 10, the test accuracy moves from 83.3\% to 84.4\%.

The train loss went down, from 0.818 at epoch 1 to 0.109 at epoch 10, and it was
almost flat after epoch 4 (0.146 at epoch 4, 0.109 at epoch 10). The test loss
did the opposite. It fell to its lowest 0.506 at epoch 2 and then went up at
every single epoch after that, to 0.747 at epoch 10. So the model kept getting
better on the 360 pictures it was trained on and worse on the 90 it had not
seen, which is what overfitting looks like.

Train 96.1\% against test 83.3\% is a gap of 12.8 points, and the other runs
show the same shape (long 95.6\% / 84.4\%, frozen 85.8\% / 74.4\%). With only
360 training images of one toy line, the model can partly remember the pictures
instead of learning the characters. The test images come from the same search
and are the same kind of product photo as the training ones, so 83.3\% is an
optimistic number: on pictures from another source the same model gets 70.9\%
(Problem 4).

\begin{center}\picturefile[0.5\linewidth]{figures/run_confusion.png}\end{center}

It mixes up Chiikawa and Hachiware most: 5 Chiikawa pictures were called
Hachiware and 2 Hachiware pictures were called Chiikawa, 7 of the 15 mistakes.
Chiikawa and Usagi come next (2 + 4 = 6), and Hachiware and Usagi are almost
never confused with each other (2). Chiikawa being in both pairs does not
surprise me: it is the plain one, a round white body with no ears and no colour
of its own, so both of the coloured characters can look like it. What did
surprise me is that the mistakes are not spread evenly: 58 of the 90 test
pictures sit on the diagonal, and the rest are all rumours of Chiikawa.

Same images, same code, same 10 epochs, and the two test accuracies are 83.3\%
with the ImageNet weights and 66.7\% from random numbers: 16.6 points apart.
The ImageNet weights have already seen a million photographs, so edges, fabric,
fur and round shapes are there before training starts; starting from random
numbers, everything has to be learned from the 360 images we have, and 10
epochs is not enough for that. The 16.6 points are the price of not borrowing
someone else's 1.2 million pictures.

Training only the last layer reaches 74.4\% against 83.3\% for all layers, 8.9
points behind, while it changes 1,539 of the 11,178,051 parameters, about
0.01\% of them. For a three-class problem the last layer is most of the
decision, and the layers below only have to move a little; the 8.9 points are
what tuning the layers below is worth here. My own setting, 20 epochs at the
same learning rate, reached 84.4\%, only 1.1 points more for twice the time.

\subsection*{Output of the runs}
\outputfile{outputs/run_output.txt}
\outputfile{outputs/scratch_output.txt}
\outputfile{outputs/frozen_output.txt}
% One more line like these for the setting you chose, named by its tag.
\outputfile{outputs/long_output.txt}

% ---------------------------------------------------------------------------
\section*{Problem 4. Cleaning}

\subsection*{My rule}
An image stays only if the main thing in it is exactly one of the three
characters, clearly, and nothing else of the story is in the frame. It goes if
the character is not the one the folder says, if it is a person or a drawing
that does not look like the character at all, or if two or more characters share
the picture.

\picturefile[0.9\linewidth]{figures/look_chiikawa_1.jpg}

\texttt{clean.py look} drew all 450 downloaded images onto 12 sheets, so that
every picture could be looked at; this is the first of the four chiikawa sheets.
I wrote the rule before starting and decided every image by the rule, not by
what looked wrong.

\subsection*{Three kinds of junk}
\sidebyside{figures/junk_1.jpg}{the folder says Chiikawa and the picture is
Usagi, drawn on a yellow background: the wrong character}
           {figures/junk_2.jpg}{not a character at all: a person, drawn, and the
watermark of the tool that drew her}

\picturefile[0.45\linewidth]{figures/junk_3.jpg}

The third kind: a shelf of goods with several characters in one frame. Whoever
made it meant the shelf, not one character.

\subsection*{What I removed}
\outputfile{outputs/clean_count.txt}

That table counts only the files I moved with \texttt{clean.py remove}, because
the log of the batch removal that followed was written with its columns in a
different order and \texttt{clean.py count} cannot read those rows: they are
counted in the \texttt{other} column. \texttt{data/removed/} holds 271 files in
all: 134 chiikawa, 38 hachiware and 99 usagi.

\picturefile[0.9\linewidth]{figures/suspect_copies.jpg}

\texttt{clean.py suspects} ran on the 450 images of 8 October, before I cleaned
anything, and reported 133 pairs of pictures closer to each other than 0.95:
that number is the title of \texttt{copies.png}, and it is the number of
\texttt{copy of} rows in \texttt{suspects.csv}. 56 of the pairs, counted from
the same rows, are two different classes, that is the same picture filed under
two characters, which cannot both be right. Counting those 133 rows against what
is in \texttt{data/clean} today, 87 pairs have neither file left, 19 pairs have
one file left and 27 pairs still have both, and those 27 are the ones to worry
about: the split into training and test is random, so one of a near-copy pair
can be in the training set while its partner is in the test set. That is one
reason the test accuracy of the cleaned model reads 100\%, and it is not a
reason to believe it. I never ran \texttt{clean.py suspects} again after adding
the 138 extra images, so the newest pictures in \texttt{data/clean} have not
been checked for near copies at all.

Of the first 20 suspects of each class, the 20 images each class agrees with its
own label least, I removed 19 of 20 for chiikawa, 14 of 20 for hachiware and 13
of 20 for usagi: 46 of 60. Looking at the sheets, the suspects did find junk I
had missed when I looked myself: wrong characters, fan drawings and pictures
with two characters in them were exactly what sat at the bottom of each class,
with their own label at 4.4\%, 5.5\% and 7.0\% probability on average. The
other way round also happened: 14 of those 60 are pictures I wanted to keep,
like a plush photographed against a busy shop shelf, which is a legitimate
picture of the character that simply does not look like the rest of the class.

\sidebyside{figures/suspects_chiikawa.jpg}{the 32 chiikawa pictures that look
least like the rest of chiikawa (rank 1 at the top left); 19 of the first 20
went}
           {figures/suspects_usagi.jpg}{the same for usagi: 13 of the first 20
are gone}

\subsection*{Before and after cleaning}
\outputfile{outputs/compare.txt}
\outputfile{outputs/clean_output.txt}
\outputfile{outputs/clean3_output.txt}

The bottom rows of \texttt{compare.txt} are the fair comparison, because the
new images are the only ones cleaning did not touch: the model trained before
cleaning gets 70.9\% on them (61 of 86) and the one trained after gets 73.3\%
(63 of 86), so cleaning was worth 2.4 points on pictures from another source.
On the downloaded test set the same rows read 83.3\% and 97.1\%, but that set
is not the same set of pictures before and after: cleaning changed it and made
it much smaller (90 images before, 35 after), so those two numbers cannot be
compared directly.

Two things about \texttt{compare.txt} that are not what they look like. The
counts of training images at the top, 104/110/99, are \texttt{data/clean} today,
not the folder the run tagged \texttt{clean} was trained on: that run read 179
images, 16 chiikawa, 112 hachiware and 51 usagi, as the first lines of
\texttt{clean\_output.txt} show, because at that moment most of the downloaded
chiikawa and usagi pictures had just been removed as junk. And \texttt{data/clean}
today holds 138 images that are not in \texttt{data/raw} (88 chiikawa, 50
usagi), added afterwards so that the two small classes would have enough
pictures to learn from; the table in Problem 4 counts the downloaded ones only.

I trained after cleaning three times, as the folder grew: \texttt{clean} on 179
images (16 chiikawa, 112 hachiware, 51 usagi), \texttt{clean2} on 208 (38, 112,
58) once more chiikawa pictures were in, and \texttt{clean3} on 315 (104, 112,
99). The runs are printed below, and only \texttt{clean3} is the one exported
for the web page.

\outputfile{outputs/clean2_output.txt}

The model I publish is the run tagged \texttt{clean3}: 315 images, 104 chiikawa,
112 hachiware, 99 usagi, 100.0\% train and 100.0\% test accuracy. The 100\% is
not believable and I do not believe it. The test set comes from the same search
as the training set, some near copies are in both halves of the split, and 315
images are few. \texttt{check.py} measures the same model on the new images and
gets 78 of 86, 90.7\%; that is the number I trust, and it is still optimistic
for the reason just given.

\sidebyside{figures/clean3_curves.png}{clean3: the model published on the web
page}
           {figures/clean3_confusion.png}{its confusion matrix on the downloaded
test set, everything on the diagonal}

% ---------------------------------------------------------------------------
\section*{Problem 5. Images from a new source}

My new images are photographs of goods my friends own. I asked them to
photograph their own plushes, figures and keyrings and send me the pictures, so
every one of them was made by a camera rather than found by the search that
filled \texttt{data/raw}: the search returned shop shots and drawings, and these
are ordinary phone pictures of toys. That is also why they cannot be among the
downloads: the model never searched, and a picture nobody uploaded to a shop is
not in a search result. \texttt{clean.py overlap} agrees: none of the 86 new
images is a near copy of any of the 571 downloaded ones, and the closest pair
is 0.95, below the 0.95 line that command uses.

\outputfile{outputs/overlap.txt}

\textbf{On the downloaded test set} the model I publish scores 100.0\% (315
images, the first lines of \texttt{clean3\_output.txt}); \textbf{on the new
images} it scores 78 of 86, 90.7\%. The nine points between them are the most
useful number in this report: the first one is measured on pictures that came
out of the same search as the training pictures, the second one on pictures that
did not.

\subsection*{Five mistakes it was confident about}

\picturefile[0.35\linewidth]{figures/mistake_1.jpg}\\
usagi\_character/이미지 3.jpg, said Chiikawa (92\%). A plush held
up in my hand, in a shop. From this angle nothing in the picture says Usagi: the
long ears that make Usagi Usagi are not in the frame, and what is left is a
round cream face with pink cheeks and a small mouth, which is what 104 Chiikawa
pictures in the training set look like.

\picturefile[0.25\linewidth]{figures/mistake_2.jpg}\\
chiikawa\_character/images-2.jpeg, said Hachiware (88\%). A Chiikawa drawing
printed very small in the middle of a large pink card. The card is nearly
square, so \texttt{prepare\_image} keeps almost all of it, and almost all of it
is flat pink: the character is a handful of thin lines low in the frame. There
was nothing in the 224x224 picture for the model to recognise.

\picturefile[0.3\linewidth]{figures/mistake_3.jpg}\\
usagi\_character/이미지 33 2.jpg, said Chiikawa (65\%). A plush
lying on dark clothing, photographed from above. The ears are out of the frame
or folded back, the background is a dark, high-contrast surface that has nothing
to do with any training picture, and the object itself is a plain cream face
with pink cheeks.

\picturefile[0.3\linewidth]{figures/mistake_4.jpg}\\
usagi\_character/이미지 29 2.jpg, said Chiikawa (57\%). A shop
pegboard: several plushes, price tags and boxes, and the Usagi is a small part
of the frame. This is the picture that is unlike the training set in the most
ways at once: a third object in a tool that only looks at the middle of the
frame, a real background, and one character among many.

\picturefile[0.3\linewidth]{figures/mistake_5.jpg}\\
usagi\_character/이미지 8.jpg, said Chiikawa (54\%). A close-up of
a plush against a rainy window, taken into the light. The object is large in the
frame, so this is not a distance problem; the light is grey, the plush is almost
colourless, and the round face with pink cheeks is all that is left to see.

Four of the five mistakes are Usagi called Chiikawa, and each one of them hides
the ears, which is the difference between the two characters. The fifth is a
Chiikawa seen so small that the crop threw it away. Both kinds say the same
thing: the model needs the character centred, large and with its ears.

\subsection*{What is different about my new images}
Every one of them is a photograph of a three-dimensional object: plushes,
figures, keyrings, taken with a phone from many angles, close up, from above,
held in a hand, indoors and outdoors, on desks, on beds, on clothes and in
shops. My downloads are mostly flat drawings and product shots: the object
centred, one colour or white behind it, even light, no hand, no shop, no
weather. The new pictures have real backgrounds, uneven light, and the object
anywhere in the frame, and they show a different line of merchandise than the
drawings do.

% ---------------------------------------------------------------------------
\section*{Problem 6. The demo}

Address: \url{https://demogom.github.io/mas1004-assignment1/}

I showed the page things that are none of my three categories. With my own face
filling the frame it answered Usagi, and with my upper body in the picture
instead, blue clothes included, it answered Hachiware. It never answers
``not a character'', because the page can only choose between the three classes:
whatever you put in front of it comes out as the class it is least unlike.

% ---------------------------------------------------------------------------
\section*{Problem 7. What the model looks at}

% Copy this subsection once for every guess.
\subsection*{Guess 1}

I think the model goes mostly by colour. When I showed the web page things that
are not one of the three characters at all, it answered Usagi for my own face,
and Hachiware as soon as my blue clothes were in the picture; anything yellow in
front of the camera came out as Usagi, and most other colours came out as
Hachiware. The three characters are told apart by colour in the drawings, too:
Usagi is yellow, Hachiware is white and blue, Chiikawa is plain white. The
mistakes point the same way, but weakly: four of the five confident mistakes in
Problem 5 are Usagi called Chiikawa, and in those pictures the yellow that says
Usagi is either missing from the object or hidden by the dark background and the
grey light.

The change I chose is to take the colour away: grayscale copies of the pictures,
with the originals untouched, so nothing else about them changes. If the model
is really reading colour, greys should cost it a lot of accuracy; if it is
reading shape, the same pictures in grey should come out almost the same. I
asked the agent to make the grayscale copies of 10 images per class and to tell
me the model's answers, so that the loss, if there is one, can be counted rather
than guessed.

My request to the agent, word for word:
\begin{quote}\small
chiikawa\_character: 0021,0029,0036,0095,3045a4591f30816d69459cf632a95fc0,0119,이미지 16 2,이미지 46,c0348d114447ffc95d74bcd4ee2e4a89,e20ee5bd663fe3c2dc4df83a13d56a7b\\
hachiware\_character: 0014,0033,0030,0047,0064,0079,0096,0120,0149,0128\\
usagi\_character: 이미지 35, 이미지 36, 이미지 33, 0020,0023,0076,0092,bbea743755f3374e8676735886c155c8, 이미지4, 이미지23\\[0.4em]
이렇게 10장을 흑백으로 해주세요 실수 5장은 그대로 두세요
\end{quote}

The change: ten pictures in each class, thirty in all, copied out of
\texttt{data/clean} into \texttt{data/changed} and turned grey there, so the
originals are untouched and the only thing that differs inside a pair is the
colour. The five confident mistakes of Problem 5 are not among the thirty: they
are in \texttt{data/new\_images}, and they are the pictures my guess was built
on, so changing them would have tested nothing new. The answers were worked out
by \texttt{src/changed\_answers.py}, which loads the same \texttt{docs/model.onnx}
the web page loads, so the sheets below, the table below them and the page all
speak from one model.

Each sheet holds ten pairs, one per line: on the left the original picture, on
the right the grey copy, and under each of them the model's answer and its
probability.

\textbf{The ten chiikawa pictures.}
\picturefile[0.8\linewidth]{figures/grey_chiikawa_character_sheet.png}

\textbf{The ten hachiware pictures.}
\picturefile[0.8\linewidth]{figures/grey_hachiware_character_sheet.png}

\textbf{The ten usagi pictures.}
\picturefile[0.8\linewidth]{figures/grey_usagi_character_sheet.png}

Read off the table: 30 of 30 right before the change, 100.0\%, and 27 of 30
after it, 90.0\%. Per class, chiikawa 10 and 10, hachiware 10 and 9, usagi 10
and 8, and all three answers that moved moved to Chiikawa. Looking down the
probabilities, the two coloured classes lose confidence when the colour goes -
in 19 of the 20 hachiware and usagi rows the after column is lower, by 1 to 21
points - while the ten chiikawa rows do not move at all: eight of them read
100\% in both columns, one goes 99\% to 98\%, and one goes up, from 94\% to
100\%.

\outputfile{outputs/changed_answers.txt}

The three answers that moved, the original next to its grey copy:

\sidebyside{figures/grey_18_original.jpg}{hachiware\_character/0120.jpg, original:
Hachiware (90\%)}
           {figures/grey_18_grey.jpg}{the same picture in grey: Chiikawa (69\%)}

\sidebyside{figures/grey_26_original.jpg}{usagi\_character/이미지 23.jpg, original:
Usagi (92\%)}
           {figures/grey_26_grey.jpg}{the same picture in grey: Chiikawa (88\%)}

\sidebyside{figures/grey_30_original.jpg}{usagi\_character/이미지 4.jpg, original:
Usagi (100\%)}
           {figures/grey_30_grey.jpg}{the same picture in grey: Chiikawa (97\%)}

\textbf{What the results say.} The guess was that the model goes mostly by
colour, and the test was to take the colour away. Taking it away cost three
answers in thirty; the model still tells 27 of the 30 apart with no colour at
all. My guess expected much more than that, so it is only partly right, and the
way it is wrong is more useful than the number. All three losses go the same
way, towards Chiikawa, the plain white character, and all three are pictures
where colour was already doing very little: 0120 is a drawing whose blue hair,
blue bow and blue sky become one grey in the copy, and 이미지 23 and 이미지 4 are
photographs of pale cream plushes, one held in a hand over a wooden floor and
one lying on a knitted mat with black lace behind it, where there is hardly any
yellow to lose in the first place. The part of the guess that survived is the
direction of the errors: colour decides for the two coloured classes and does
nothing for Chiikawa. That is the same direction as the confusion matrix in
Problem 3 and the same direction as the five mistakes in Problem 5, where four
of the five are coloured characters called Chiikawa. The model uses the colour
of the character as a hint, but the hint is not what keeps it right on 90\% of
pictures that have no colour left.

\subsection*{What I now think the model uses}

I now think the model goes by the shape it can see in the middle of the frame,
and above all by the ears, with the character's own colour as a weaker second
signal that usually agrees with the shape and settles the pictures where the
shape is hard to see. Four of the five confident mistakes in Problem 5 and all
three answers that moved here are the same story: the long ears are out of the
frame, folded back, or the whole character is a small part of a busy picture, and
what is left is a round cream face with pink cheeks, which is exactly Chiikawa.
The one mistake that goes the other way, a Chiikawa drawing printed so small
that the crop throws it away, is the same story told from the other side: too
little of the character in the 224x224 square to see anything. Colour is real
but secondary, and the grey copies show how secondary: the ten Chiikawa rows
hardly move at all, because that character has no colour to lose, while the
coloured rows only lose confidence, not answers.

How sure am I? Fairly sure about the ears, and sure that Chiikawa is the class
everything falls back to, because in every mistake I have looked at Chiikawa is
on one side or the other. Less sure about colour, for two reasons. Thirty
pictures, ten for each class, is not many, and my change was blunt: it takes
colour out of the character and the background at once, so it cannot say which of
the two the model was reading. If I had another run, I would change the character
only and leave the background alone, and for Usagi I would take the yellow off
the character by itself, to see whether the confidence that fell in the grey rows
was the yellow of the rabbit or the colour around it.

% ---------------------------------------------------------------------------
\section*{Appendix. Output of check.py}
\outputfile{outputs/check.txt}

\end{document}
